\documentclass[10pt, twocolumn, landscape, a4paper]{article}

\usepackage[margin=1.25cm]{geometry}
\usepackage[utf8]{inputenc}
\usepackage[spanish,es-noshorthands]{babel}
\usepackage{amssymb, amsmath, amsbsy}
\usepackage{verbatim}
\usepackage{mathpazo}
\usepackage{tasks}
\usepackage{enumerate}
\usepackage[pdftex]{hyperref}
\hypersetup{pdftitle={},pdfauthor={Luis Carrillo Gutiérrez}}

\renewcommand{\familydefault}{\sfdefault}
\pagestyle{empty} 

\begin{document}
\subsection*{Biología}
\begin{itemize}
\item{Muchas \emph{bacterias patógenas} tienen una estructura gruesa de \textbf{polisacáridos gelatinosos}, que toma el nombre de :
\begin{tasks}(3)
\task{Cápside}\task{Vaina}\task{Cápsula}\task{Capsómero}\task{Pared celular}
\end{tasks}
}
\item{El \emph{electrolito cloro} para ingresar en la célula utiliza a la proteína transportadora \\ del tipo : 
\begin{tasks}(3)
\task{Uniporte}\task{Monotópica}\task{Simporte}\task{Permeasa}\task{Antiporte}
\end{tasks}
}
\item{En la membrama del \emph{retículo endoplasmático liso}, se encuentran enzimas incrustadas que catalizan la síntesis de :
\begin{tasks}(3)
\task{Glucoproteínas}\task{Fosfolípidos}\task{Gonadotropinas}\task{Peroxidasas}\task{Sustancias pécticas}
\end{tasks}
}
\item{En las \emph{células germinales}, la enzima activa que continuamente agrega DNA a los extremos de los cromosomas con lo que se evita su acortamiento, se denomina :
\begin{tasks}(3)
\task{Acrosina}\task{Nucleasa}\task{Separasa}\task{Kinasa}\task{Telomerasa}
\end{tasks}
}
\end{itemize}
\end{document}
